\section{Use case 5: Traffic flow modeling}
\label{sec:traffic-flow}

\paragraph{Problem statement}
Macroscopic traffic models describe the evolution of vehicle density
$\rho(x,t)$ along a road. In the Lighthill--Whitham--Richards (LWR) model,
density satisfies the scalar conservation law
\begin{equation}
    \partial_t \rho + \partial_x q(\rho) = 0,
    \label{eq:traffic-lwr}
\end{equation}
where the fundamental diagram $q(\rho)$ relates density to traffic flux and
the corresponding velocity is $v(\rho)=q(\rho)/\rho$. Classical fundamental
diagrams provide compact and physically interpretable constitutive laws, but a
fixed analytical form can retain systematic bias when it is embedded in a
spatio-temporal traffic simulation. Replacing the diagram entirely with a
flexible data-driven model may improve fit at the cost of interpretability and
physical structure.

We therefore considered the intermediate hypothesis that a calibrated
classical diagram captures the principal traffic physics, while its remaining
bias can be represented by a compact multiplicative correction,
\begin{equation}
    q_{\mathrm{corr}}(\rho)
    = q_{\mathrm{base}}(\rho;\boldsymbol{\theta})
      g(\rho;\boldsymbol{c}).
    \label{eq:traffic-correction}
\end{equation}
Here $q_{\mathrm{base}}$ is one of five established diagrams---Greenshields,
the equilibrium Intelligent Driver Model (IDM), Weidmann, Triangular, or Del
Castillo---whose calibrated coefficients $\boldsymbol{\theta}$ remain fixed
during structural discovery. The correction $g$ is dimensionless, and
$g=1$ recovers the unmodified diagram. Restricting the selected corrections to
positive exponential multipliers preserves the units, sign, and zeros of the
baseline flux, while still allowing local and short-range spatial adjustments.

\paragraph{Method}
The \texttt{automodel} skill was used with OpenAI Codex with the following
domain-specific choices on top of the baseline workflow.

\noindent\textbf{(i) Objective and data splitting.}
We used density, velocity, and flow measurements from the NGSIM Interstate~80
data set. After nondimensionalization, the data comprise 79 spatial rows and
180 time samples; 75 interior spatial positions enter the objective. The split
was chronological and fixed before model discovery: samples 0--63 were used to
fit correction coefficients, samples 64--107 for structural validation and
model selection, and samples 108--179 as a final held-out test interval. No
shuffling was used. The two component errors are relative $L_2$ errors,
\begin{equation}
    E_\rho = \frac{\lVert\widehat{\rho}-\rho\rVert_2}
                        {\lVert\rho\rVert_2},
    \qquad
    E_v = \frac{\lVert\widehat{v}-v\rVert_2}
                    {\lVert v\rVert_2},
\end{equation}
and candidate fit was summarized by
\begin{equation}
    E_{\mathrm{data}} = 50\left(E_\rho^2+E_v^2\right).
    \label{eq:traffic-data-error}
\end{equation}
Structure selection used validation
$E_{\mathrm{fitness}}=E_{\mathrm{data}}+0.01n$, where $n$ is the number of
nodes in the expression tree. Thus, an added mechanism had to offset its
complexity cost rather than merely produce a numerically negligible reduction
in residual error.

\noindent\textbf{(ii) Model and evaluator.}
A candidate model was a multiplier $g$ applied to one of the five fixed
fundamental diagrams in Eq.~\eqref{eq:traffic-correction}. Each candidate was
evaluated by solving Eq.~\eqref{eq:traffic-lwr} with the same initial state,
measured boundary density, mesh, and time refinement as its baseline. The
finite-volume update used a Rusanov/Godunov flux implemented with discrete
exterior calculus operators. For nonlocal corrections, the wave-speed bound
included the off-diagonal terms of the flux Jacobian, rather than treating a
spatially coupled flux as pointwise.

The coefficient optimizer used deterministic bounded Powell searches with
explicit starts and evaluation budgets. Candidates were rejected if their
multipliers or simulations were non-finite, if they produced a non-positive
multiplier or negative velocity on the relevant physical density domain, or if
the corrected velocity violated the prescribed monotonicity check. Every
expression, parameter vector, train/validation score, feasibility result,
runtime, and memory measurement was retained.

\noindent\textbf{(iii) Structural search and confirmation.}
The search used $M\times S\times I=5\times3\times3$: five meta-iterations,
three parallel hypothesis streams per iteration, and three sequential
structural attempts per stream. Every proposed structure was fitted separately
to all five baseline diagrams, giving 225 baseline--candidate fits before final
refitting. Early iterations explored positive pointwise polynomial, anchored,
saturating, and rational corrections. Later iterations introduced short-range
convolutions and Hodge-star features defined on the traffic mesh, and refined
the most promising incumbents. Shortlisted candidates were re-evaluated in the
root process before acceptance, and the test interval was not accessed during
this search.

Let
\begin{equation}
    S_3(\rho)=\operatorname{conv}_3(\rho,\mathbf{1}),
    \qquad
    C(\rho)=S_3(\rho)-3\operatorname{conv}_1(\rho,\mathbf{1}),
\end{equation}
where $C$ is a local density-contrast feature, and let $H_2(\rho)$ denote the
quadratic Hodge-star feature obtained by squaring the dual one-cochain
$\star\rho$ and mapping it back to primal nodes. Validation selected the
following baseline-specific structures:
\begin{align}
    g_{\mathrm{G}}(\rho)
        &= \exp\!\left(c_0+a\rho+b\rho^2+dC(\rho)\right), \\
    g_{\mathrm{IDM}}(\rho)
        &= \exp\!\left(aC(\rho)+bH_2(\rho)C(\rho)\right), \\
    g_{\mathrm{W}}(\rho)
        &= \exp\!\left(a\rho\left(1-\rho/\rho_{\max}\right)\right), \\
    g_{\mathrm{T}}(\rho)
        &= \exp\!\left(a\rho C(\rho)\right), \\
    g_{\mathrm{DC}}(\rho)
        &= \exp\!\left(c_0+aS_3(\rho)+bC(\rho)\right).
    \label{eq:traffic-selected-corrections}
\end{align}
After these structures were frozen, their coefficients were refitted on the
combined train and validation interval (samples 0--107). Each frozen model was
then evaluated once on samples 108--179. The held-out result could not affect
either the selected structure or its coefficients.

As an aggregate traffic-level diagnostic, we also computed total travel time
over the held-out interval,
\begin{equation}
    \mathrm{TTS}=\int_{t\in\mathcal{T}_{\mathrm{test}}}
                 \int_{x\in\Omega}\rho(x,t)\,\mathrm{d}x\,\mathrm{d}t,
    \qquad
    E_{\mathrm{TTS}}
      =\frac{|\mathrm{TTS}_{\mathrm{model}}-\mathrm{TTS}_{\mathrm{data}}|}
             {\mathrm{TTS}_{\mathrm{data}}}.
    \label{eq:traffic-tts}
\end{equation}
Models were ranked only by held-out $E_{\mathrm{data}}$. Improvement relative
to a model's own classical baseline was defined as
$100(E_{\mathrm{data}}^{\mathrm{base}}-E_{\mathrm{data}}^{\mathrm{model}})
/E_{\mathrm{data}}^{\mathrm{base}}$.

\paragraph{Results}
Table~\ref{tab:traffic-test-results} compares the classical diagrams, the
corresponding symbolic-regression (SR) corrections, and the structures found in
the Automodel search. SR-IDM attained the lowest held-out $E_{\mathrm{data}}$
(6.226), followed by SR-Triangular (6.728), with Automodel--Del Castillo third
(6.780). SR-Del Castillo achieved the lowest density and TTS errors. The
Automodel--Greenshields correction produced the largest improvement over its
own baseline, reducing test $E_{\mathrm{data}}$ by 23.17\%, and also reduced
$E_{\mathrm{TTS}}$ from 0.152 to 0.121. Automodel improved held-out
$E_{\mathrm{data}}$ for four of the five baseline families; the exception was
IDM, for which the score increased by 2.30\%. The TTS diagnostic was not part of
model selection and did not improve consistently across the selected
corrections.

\begin{table}[t]
    \centering
    \small
    \setlength{\tabcolsep}{3.5pt}
    \caption{Held-out I80 prediction errors. Lower is better for all error
    columns. Parentheses give rank by full-precision test
    $E_{\mathrm{data}}$ across all 15 models. $\Delta_{\mathrm{base}}$ is the
    percentage reduction in test $E_{\mathrm{data}}$ relative to the model's
    own classical baseline; positive values indicate improvement. AM denotes
    \texttt{automodel}.}
    \label{tab:traffic-test-results}
    \begin{tabular}{@{}lrrrrr@{}}
        \toprule
        Model & $E_\rho$ & $E_v$ & $E_{\mathrm{data}}$ (rank)
              & $\Delta_{\mathrm{base}}$ & $E_{\mathrm{TTS}}$ \\
        \midrule
        Greenshields       & 0.294 & 0.313 & 9.230 (15) & --       & 0.152 \\
        SR--Greenshields   & 0.293 & 0.294 & 8.603 (14) & $+6.80\%$  & 0.163 \\
        AM--Greenshields   & 0.255 & 0.277 & 7.092 (7)  & $+23.17\%$ & 0.121 \\
        IDM                & 0.258 & 0.270 & 6.957 (6)  & --       & 0.116 \\
        SR--IDM            & 0.246 & \textbf{0.253} & \textbf{6.226 (1)}
                           & $+5.26\%$ & 0.108 \\
        AM--IDM            & 0.256 & 0.277 & 7.117 (8)  & $-2.30\%$ & 0.125 \\
        Weidmann           & 0.265 & 0.288 & 7.651 (12) & --       & 0.115 \\
        SR--Weidmann       & 0.255 & 0.269 & 6.861 (5)  & $+10.33\%$ & 0.121 \\
        AM--Weidmann       & 0.265 & 0.281 & 7.459 (11) & $+2.51\%$ & 0.117 \\
        Triangular         & 0.269 & 0.294 & 7.919 (13) & --       & 0.129 \\
        SR--Triangular     & 0.252 & 0.266 & 6.728 (2)  & $+15.04\%$ & 0.121 \\
        AM--Triangular     & 0.256 & 0.277 & 7.118 (9)  & $+10.11\%$ & 0.127 \\
        Del Castillo       & 0.257 & 0.279 & 7.197 (10) & --       & 0.119 \\
        SR--Del Castillo   & \textbf{0.244} & 0.276 & 6.801 (4)
                           & $+5.51\%$ & \textbf{0.107} \\
        AM--Del Castillo   & 0.253 & 0.267 & 6.780 (3)  & $+5.80\%$ & 0.123 \\
        \bottomrule
    \end{tabular}
\end{table}

The SR values in Table~\ref{tab:traffic-test-results} were reproduced with the
official SR-Traffic results implementation at commit \texttt{3bf285a}, whereas
the classical and Automodel rows were produced by the present checkout. Each SR
improvement is paired with the classical result from that same upstream
implementation; in particular, its IDM baseline has
$E_{\mathrm{data}}=6.571$. The rankings are therefore a reproduced benchmark
rather than a claim that all models share an identical implementation path.
Moreover, this study uses one road and one chronological split. The simulation
is conditional on measured boundary density throughout the test interval, and
the benchmark normalization uses the maximum density over the complete time
series. These choices preserve compatibility with the repository benchmark but
limit the strength of claims about fully autonomous forecasting and
out-of-distribution generalization. The results support multiplicative
correction as a useful baseline-dependent strategy, but do not show that
Automodel uniformly outperforms the existing SR models.
